% ---
% id: meat-core/rem:tagged-entries-and-accumulation
% kind: rem
% title: "Party-local entries and constrained accumulation"
% summary: >-
%   No statement of the constrained accumulation theorem becomes false for a
%   list with a party-local entry: its hypotheses are satisfiable, its
%   conclusion holds and its monotonicity survives. The tagged reading is not
%   shown canonical, and the paragraph before the theorem describes only
%   untagged forms.
% prominence: standard
% status: accepted
% uses: [protocol-compiler/def:homog, thm:constrained-accumulation, def:roundmodel, lem:range, def:hqkd-list]
% mentions: [rem:homog, protocol-compiler/rem:party-local-soundness]
% assumes: []
% refs: []
% provenance: newly-derived
% lean:
% ---

\begin{remark}\label{rem:tagged-entries-and-accumulation}
Theorem~\ref{thm:constrained-accumulation}, which turns a licensed list into a
key-length licence, is stated for a list licensed for the declared mode and
for an \tref{attack sequence} admissible for that list. Both readings carry
the \tref{party-local form} of \xref{protocol-compiler/def:homog}, the
licensed class being that definition's and the per-round admissibility format
of Definition~\ref{def:roundmodel} carrying the party-local admissibility set
beside the two of the untagged forms, with
Lemma~\ref{lem:range}\,(i) transporting it to the state level in either
mode. \emph{So no statement of that theorem becomes false of a list carrying
a tagged entry}: its hypotheses are satisfiable, its conclusion holds by that
transport, and its monotonicity $\hqkd(\CList)\ge\hqkd(\varnothing)$, where
$\hqkd(\CList)$ is that theorem's single-round QKD rate at order $\alpha$
under the constraints of $\CList$, holds because a tagged entry shrinks the
\tref{feasible set} further than its untagged reading does.

\emph{What it does not say} is that the tagged reading is canonical: the
argument that the licensed class is not an arbitrary choice is stated for the
two untagged forms (Remark~\ref{rem:homog}) and is not extended to the third
here, and the third form's soundness rests on the certificate on the scoped
factor's own channel (\xref{protocol-compiler/rem:party-local-soundness}) and
on nothing this contract proves.

\emph{The paragraph motivating the theorem describes the untagged forms.} That
paragraph, before the theorem's statement, rests it on three points: the
imported accumulation theorem is stated per sequence of attack channels;
which attack channels the adversary is granted is a modelling input outside
that theorem; and a constraint list licensed for the declared mode is
``precisely the state-level shadow of an attack-set restriction in the
published format'', the homogeneous entries of form (a) and the scalar
entries of form (b) of \xref{protocol-compiler/def:homog} each being carried
to the state level by a transport of its own. \emph{The third point is stated
of the untagged forms in its main clause, not only in the two transports it
lists}: the published format it cites is the single comparison of a channel
adjoint against a multiple of the identity, whereas the admissibility set of
a tagged entry is a \emph{family} of operator inequalities, one per reference
system and per positive semidefinite operator on it
(Definition~\ref{def:roundmodel}). \emph{What covers the third form is the
second point}, which makes restricting attention to a declared subset of
attack channels a modelling step and is stated for any declared subset, not
for a shape of right-hand side. So the theorem's hypotheses, conclusion and
proof are unaffected, and only the description in its motivating paragraph
does not reach the third form.
\end{remark}
